%==============================================================================
% فصل دوم: بررسی متون
% دو بخش اصلی: (۱) مبانی نظری و مفهومی  (۲) مرور مطالعات مرتبط
%==============================================================================
\ChapterStart
\chapter{بررسی متون}
\label{ch:literature}


\section{مبانی نظری و مفهومی}
\label{sec:theoretical}

بیماری‌های قلبی عروقی، به ویژه انفارکتوس میوکارد (\lr{MI})، یکی از عوامل اصلی مرگ و میر و ناتوانی در سراسر جهان و همچنین در ایران به شمار می‌روند و بار اقتصادی و بهداشتی سنگینی را بر جوامع تحمیل می‌کنند\cite{GBD2021Causes2024,HeidariForoozan2025Iran}. با وجود پیشرفت‌در درمان‌ حاد، بیماران پس از ترخیص همچنان در معرض خطر بالای عود حوادث ایسکمیک قرار دارند. به همین دلیل، پیشگیری ثانویه شامل مداخلات دارویی و به خصوص اصلاحات غیردارویی سبک زندگی نظیر ترک دخانیات، فعالیت بدنی منظم، رژیم غذایی سالم و مدیریت استرس، نقش بسیار حیاتی در کاهش خطر مرگ و میر و عوارض بعدی ایفا می‌کند \cite{Byrne2023ESC}. اجرای دقیق این دستورالعمل‌ها می‌تواند کیفیت زندگی بیماران را به میزان قابل توجهی ارتقا داده و از بروز مجدد بیماری جلوگیری نماید\cite{Dibben2021CardiacRehabilitation}.

با وجود اهمیت اثبات‌شده پیشگیری ثانویه، شواهد بالینی نشان می‌دهد که میزان پیروی بیماران از توصیه‌های اصلاح سبک زندگی و خود‌مراقبتی در طول زمان به شدت کاهش می‌یابد. عدم پایبندی به این توصیه‌ها احتمال بستری مجدد و مرگ را افزایش می‌دهد. محققان دریافته‌اند که تغییر رفتار و اصلاح سبک زندگی یک فرآیند پیچیده و چندوجهی است که صرفاً با ارائه اطلاعات پزشکی به بیمار محقق نمی‌شود. از این رو، برای درک عمیق‌تر موانع و تسهیل‌کننده‌های پیروی درمانی، استفاده از تئوری‌ها و مدل‌های روان‌شناختی و رفتاری در متون علمی به عنوان یک ضرورت اجتناب‌ناپذیر مطرح شده است.

یکی از چارچوب‌های نظری برجسته در این زمینه، «مدل اعتقاد بهداشتی»\footnote{\lr{Health Belief Model}} است. این مدل فرض می‌کند که تصمیم‌گیری یک فرد برای اتخاذ یک رفتار پیشگیرانه، مستقیماً به ادراک او از تهدید بیماری و ارزیابی منافع در برابر موانع بستگی دارد. در بیماران مبتلا به انفارکتوس میوکارد، اگر بیمار خود را مستعد بروز مجدد سکته بداند (حساسیت درک‌شده) و عواقب آن را جدی و مرگبار تلقی کند (شدت درک‌شده)، تمایل بیشتری برای تغییر رفتار خواهد داشت. همچنین، بیمار باید باور داشته باشد که پیروی از رژیم غذایی کم‌نمک یا ورزش (منافع درک‌شده) بر موانعی چون کمبود وقت یا خستگی (موانع درک‌شده) غلبه دارد تا بتواند رفتار خود را با موفقیت اصلاح کند.

تئوری دیگری که در تبیین پایبندی به سبک زندگی نقش کلیدی دارد، «تئوری رفتار برنامه‌ریزی‌شده»\footnote{\lr{Theory of Planned Behavior}} است. بر اساس این نظریه، مهم‌ترین پیش‌بینی‌کننده رفتار، قصد یا نیت فرد است که خود تحت تأثیر سه عامل اصلی قرار دارد: نگرش نسبت به رفتار، هنجارهای ذهنی و کنترل رفتاری درک‌شده. در زمینه پیشگیری ثانویه، نگرش مثبت بیمار به ترک سیگار، حمایت و فشارهای مثبت از سوی خانواده و کادر درمان (هنجارهای ذهنی) و همچنین باور به توانایی خود در غلبه بر چالش‌های رژیم غذایی و مدیریت استرس (کنترل رفتاری درک‌شده)، مجموعاً قصد بیمار را برای اصلاح سبک زندگی مخرب تقویت کرده و به تبعیت پایدارتر منجر می‌شوند.

در همین راستا، «مدل ارتقای سلامت پندر»\footnote{\lr{Pender's Health Promotion Model}} نگاهی جامع‌تر به عوامل تعیین‌کننده رفتارهای بهداشتی دارد. این مدل بر خلاف رویکردهای مبتنی بر ترس از بیماری، بر تمایل مثبت فرد به سوی بهزیستی و ارتقای سلامت تأکید می‌ورزد. در این مدل، تجربیات قبلی فرد، ویژگی‌های فردی، و شناخت‌ها و عواطف مرتبط با رفتار نقش محوری دارند. پژوهش‌های مبتنی بر مدل پندر در بیماران سکته قلبی نشان داده‌اند که خودکارآمدی (اعتماد به نفس بیمار برای انجام موفقیت‌آمیز فعالیت بدنی یا مدیریت استرس در شرایط دشوار) یکی از قوی‌ترین محرک‌ها برای پایبندی طولانی‌مدت به تغییرات غیردارویی است.

در نهایت، «مدل خودتنظیمی لِوِنتال»\footnote{\lr{Leventhal's Self-Regulation Model}} با تمرکز بر مفهوم ادراک از بیماری\footnote{\lr{Illness Perception}}، بعد دیگری از این پدیده را روشن می‌سازد. طبق این مدل، بیماران بر اساس اطلاعات و تجربیات خود، بازنمایی‌های شناختی و احساسی خاصی از بیماری خود (نظیر علت، پیامدها، طول مدت و قابلیت کنترل آن) می‌سازند. بیمارانی که ماهیت مزمن بیماری عروق کرونر را درک نکرده و سکته قلبی را صرفاً یک رویداد گذرا می‌پندارند، یا کنترل پایینی بر شرایط خود احساس می‌کنند، دچار اضطراب و ترس از عود شده و راهبردهای مقابله‌ای ناکارآمدی اتخاذ می‌کنند که نتیجه آن عدم پیروی از سبک زندگی سالم است. بنابراین، اصلاح ادراک بیمار از بیماری، زیربنای اساسی برای تضمین پایبندی به اقدامات پیشگیری ثانویه محسوب می‌گردد.



\section{مروری بر مطالعات پیشین}
\label{sec:previous-studies}


بررسی پیشینه علمی در این مطالعه با استفاده از کلیدواژه‌های فارسی و انگلیسی در پایگاه‌های داده معتبر علمی داخلی و خارجی انجام پذیرفت. کلیدواژه‌های فارسی شامل "انفارکتوس میوکارد"، "سکته قلبی"، "پیشگیری ثانویه"، "اقدامات غیردارویی"، "پیروی درمانی"، "پایبندی به درمان"، و "سبک زندگی" و کلیدواژه‌های انگلیسی شامل \lr{Myocardial Infarction}»، \lr{"Secondary Prevention""}، \lr{"Non-pharmacological"}، \lr{"Adherence"}، \lr{"Compliance"}، و \lr{"Lifestyle"} بودند. جستجوی فارسی در پایگاه داده \lr{SID} و جستجوی انگلیسی در پایگاه‌های داده \lr{PubMed}، \lr{Scopus} و \lr{WoS} انجام پذیرفت. در نهایت از میان مطالعات مورد بررسی، چهار مطالعه داخلی و چهار مطالعه خارجی انتخاب گردیدند.


\subsection{مطالعات داخلی}


\begin{itemize}
\item تقی‌زاده و همکاران در یک مطالعه مقطعی توصیفی در سال ۲۰۲۲ میلادی، نقش ادراک از بیماری و چارچوب‌های شناختی را بر پذیرش بیماری و بروز عوامل خطر قلبی عروقی در تبریز ارزیابی کردند. این مطالعه بر روی ۱۳۱ بیمار که برای نخستین بار تحت مداخله آنژیوپلاستی عروق کرونر در یک بیمارستان فوق تخصصی قرار گرفته بودند انجام شد و رفتارهای مرتبط با سبک زندگی، شاخص‌های بیومتریک و درجات اضطراب و افسردگی مورد سنجش قرار گرفت. نتایج مطالعه نشان داد که ۶۶.۴ درصد از بیماران دارای حداقل یک عامل خطر زمینه‌ای بودند و ۷۸.۶ درصد از آنان مطلقاً هیچ‌گونه فعالیت ورزشی نداشتند. از نظر روانی نیز ۴۱.۲۲ درصد درگیر افسردگی و ۳۵.۸۸ درصد درگیر اضطراب بالینی بودند. ارزیابی شناختی مشخص کرد که میانگین نمره کل ادراک از بیماری ۴۳.۴۹ از ۸۰ بود و پایین‌ترین نمره بیماران مربوط به بُعد درک و شناخت از ماهیت بیماری با نمره ۳.۷۷ گزارش شد. نتایج حاکی از وجود همبستگی معناداری میان نمره کل ادراک از بیماری و میزان روزهای فعالیت بدنی (\lr{r=-0.21, p=0.014})، مصرف غذاهای چرب (\lr{r=-0.17, p=0.048}) و میزان خواب (\lr{r=-0.24, p=0.005}) بود. همچنین، همبستگی مستقیم و معناداری بین ادراک از بیماری و افسردگی (\lr{p<0.001})، اضطراب (\lr{p=0.001}) و شاخص‌های بیومتریک نظیر کلسترول (\lr{r=0.23, p=0.006}) و فشار خون سیستولیک (\lr{r=0.37, p<0.001}) اثبات گردید. محققان نتیجه گرفتند که مشکل عدم پیروی از درمان در بیماران ایرانی ریشه در نقص اطلاعاتی دارد و ادراک ضعیف از بیماری می‌تواند منجر به پیامدهای روان‌شناختی و شکست در پایبندی به سبک زندگی سالم شود. \cite{Thagizadeh2022IllnessPerception}

\item امینی و همکاران در یک مطالعه شبه‌تجربی در سال ۲۰۲۱ میلادی، اثر برنامه‌های آموزش خودمراقبتی را بر فاکتورهای خطر قلبی و مدیریت فردی سبک زندگی پس از ترخیص بیماران در استان همدان بررسی کردند. این مطالعه بر روی ۹۲ بیمار مبتلا به سکته قلبی در دو گروه آزمایش و کنترل انجام شد و گروه مداخله تحت یک برنامه فشرده هشت هفته‌ای برای پایش کیفیت رژیم غذایی، فعالیت فیزیکی در منزل از طریق شمارش گام‌ها، نوسانات فشار خون، اندازه دور کمر و شاخص توده بدنی قرار گرفتند. نتایج مطالعه نشان داد که در پایان هشت هفته، گروه آزمایش نمرات به مراتب بالاتری در رعایت رژیم غذایی نسبت به گروه کنترل کسب کرد (\lr{p<0.001}) و میزان فعالیت بدنی خانگی با افزایش محسوس در شمارش گام‌های روزانه (۲۲۱۲ در برابر ۲۱۰۶ گام) بهبود یافت (\lr{p=0.004}). با این وجود، این مداخله کوتاه‌مدت نتوانست تغییرات آماری معنی‌داری در شاخص‌های مورفولوژیک و همودینامیک با ثبات از قبیل فشار خون سیستولیک و دیاستولیک، دور کمر و شاخص توده بدنی ایجاد نماید (\lr{p>0.05}). محققان نتیجه گرفتند که برنامه‌های خودمراقبتی در کوتاه‌مدت می‌توانند عادات رفتاری مثل گام برداشتن و تغذیه را اصلاح کنند، اما برای مشاهده کاهش در فاکتورهای فیزیکی نظیر فشار خون و وزن نیاز به استمرار حمایت‌ها در بازه‌های طولانی‌تری است. \cite{Amini2021SelfManagement}

\item نجفی و همکاران در یک مطالعه هم‌گروهی گذشته‌نگر در سال ۲۰۲۴ میلادی، نابرابری‌های جنسیتی را در سطح مشارکت در بازتوانی قلبی و پیامدهای تغییر سبک زندگی طی یک بازه زمانی ۱۸ ساله در یک مرکز تخصصی در ایران بررسی کردند. این پژوهش بر روی ۲۰۶۴ بیمار شامل ۱۵۲۱ مرد و ۵۴۳ زن انجام گرفت و متغیرهای کلیدی از جمله معادل متابولیک برای آمادگی جسمانی، تغییرات شاخص توده بدنی، ترک سیگار و شیوع بیماری‌های همراه نظیر پرفشاری خون در بدو ورود به جلسات بررسی گردید. نتایج مطالعه نشان داد که شکاف جنسیتی عمیقی وجود دارد و مردان با ۷۳.۶۹ درصد بخش عمده متقاضیان را تشکیل دادند که میانگین جلسات شرکت شده توسط آنان نیز به طور معنی‌داری بیشتر بود (\lr{p<0.05}). زنان در زمان ارجاع وضعیت بالینی نامطلوب‌تری داشتند و شیوع پرفشاری خون در آنان به طور قابل‌توجهی بالاتر بود (۴۴.۹ درصد در مقابل ۲۷.۰ درصد مردان، \lr{p<0.001}). با این حال، زنانی که وارد چرخه بازتوانی شدند در ترک دخانیات بسیار موفق‌تر از مردان عمل کردند (نرخ موفقیت ۶۶.۶ درصد در برابر ۳۸.۴ درصد، \lr{p<0.001})، هرچند مردان در حوزه‌های فیزیکی پیشرفت بهتری داشتند. محققان نتیجه گرفتند که به دلیل موانع فرهنگی و لجستیکی برای زنان، پزشکان باید برنامه‌های ورزشی خانگی را برای آنان در نظر بگیرند و ارجاع مردان به مشاوره‌های روان‌شناختی جهت ترک سیگار را جدی‌تر پیگیری کنند. \cite{Najafi2024Gender}

\item آرمان و همکاران در سال ۲۰۲۵ در یک مطالعه نیمه‌تجربی به بررسی اثربخشی یک مداخله آموزشی خانواده‌محور بر پایه تئوری رفتار برنامه‌ریزی‌شده (\lr{TPB}) جهت ارتقای رفتارهای خودمراقبتی و اصلاح سبک زندگی در ۱۴۰ بیمار پس از انفارکتوس حاد میوکارد (\lr{MI}) در شهر فسا در ایران پرداختند. بیماران به صورت تصادفی به دو گروه مداخله (۷۰ نفر) و کنترل (۷۰ نفر) تقسیم شدند. گروه مداخله یک برنامه آموزشی ۱۰ جلسه‌ای را با حضور حداقل یکی از اعضای خانواده و با تمرکز بر اقدامات غیردارویی پیشگیری ثانویه، از جمله فعالیت بدنی، تغذیه سالم، مدیریت استرس و کنترل عوامل خطر دریافت کردند. ارزیابی‌های اولیه نشان داد که هیچ تفاوت معناداری در متغیرهای پایه‌ای دموگرافیک و شاخص‌های رفتاری بین دو گروه وجود نداشت (\lr{0.05<p}). با این حال، نتایج پیگیری چهار ماهه حاکی از آن بود که گروه مداخله در مقایسه با گروه کنترل، بهبود چشمگیر و از نظر آماری کاملاً معناداری را در تمامی متغیرهای ارزیابی‌شده از جمله نگرش، قصد رفتاری، نمرات سبک زندگی و رفتارهای خودمراقبتی تجربه کردند (\lr{0.001>p}). همچنین، اندازه اثر (\lr{Cohen's d}) برای این تغییرات بسیار بزرگ گزارش شد؛ به طوری که این شاخص برای نمره سبک زندگی ۴.۵۷ و برای رفتار خودمراقبتی ۱۴.۵۵ به دست آمد که نشان‌دهنده تاثیر عمیق و عملی مداخله است. محققان در این بررسی بر اهمیت نقش بافت فرهنگی ایران تاکید کرده و نشان دادند که رفتارهای خودمراقبتی و تغییرات سبک زندگی در بیماران ایرانی به شدت تحت تاثیر هنجارهای فرهنگی و حمایت‌های شبکه‌ خانواده قرار دارد. این یافته‌ها به روشنی اثبات می‌کند که برای افزایش نرخ پایبندی بیماران قلبی به توصیه‌های غیردارویی پیشگیری ثانویه، سیستم‌های بهداشتی نیازمند طراحی برنامه‌های آموزشی ساختاریافته‌ای هستند که علاوه بر هدف قرار دادن باورهای روانی-اجتماعی بیمار، خانواده را نیز به عنوان یک بازوی حمایتی در روند درمان و تغییر سبک زندگی درگیر کنند.
\end{itemize}




\subsection{مطالعات خارجی}


\begin{itemize}
\item کاتسوا و همکاران در مطالعه بزرگ چند ملیتی (\lr{EUROASPIRE V}) در سال 2019 به بررسی وضعیت کنترل ریسک‌فاکتورها و پایبندی به سبک زندگی در 8261 بیمار مبتلا به بیماری عروق کرونر در 27 کشور اروپایی و آسیایی پرداختند. نتایج این مطالعه مقطعی نشان داد که علیرغم پوشش بالای دارویی، پایبندی به اصلاحات سبک زندگی بسیار نامطلوب است. طبق یافته‌ها، 55\% از بیماران همچنان چاق بودند، 19\% سیگار می‌کشیدند و 66\% فعالیت بدنی کافی (طبق معیار 150 دقیقه در هفته) نداشتند. این مطالعه نشان داد که بین «تجویز دارو» و «اصلاح واقعی رفتار» فاصله عمیقی وجود دارد و مداخلات فعلی بازتوانی قلبی در ایجاد تغییرات پایدار در رژیم غذایی و ورزش ناکام بوده‌اند. محققان نتیجه گرفتند که سیستم‌های بهداشتی باید از رویکرد دارو-محور به سمت رویکردی جامع‌نگر با تمرکز بر مدیریت وزن و فعالیت بدنی حرکت کنند.

\item هامر و همکاران در یک مطالعه مشاهده ای و مبتنی بر داده‌های ملی در سال ۲۰۲۴ میلادی، با هدف بررسی میزان پایبندی بیماران پس از سکته قلبی به اقدامات پیشگیری ثانویه و تأثیر آن بر پیامدهای بالینی در کشور اتریش این پژوهش را انجام دادند. اهمیت این موضوع از آنجا ناشی می‌شود که علی‌رغم پیشرفت‌های درمانی، میزان بروز مجدد حوادث قلبی همچنان بالا بوده و عدم پایبندی بیماران به توصیه‌های درمانی و سبک زندگی (که از طریق برنامه‌های بازتوانی قلبی پیگیری می‌شود) یکی از عوامل کلیدی در این زمینه است. این مطالعه شامل بیماران مبتلا به انفارکتوس حاد میوکارد بود که بین سال‌های ۲۰۱۱ تا ۲۰۱۵ در اتریش ترخیص شده بودند. داده‌ها از سیستم بیمه سلامت ملی و نسخه‌های داروخانه‌ای استخراج شد و حجم نمونه برای ارزیابی مشارکت در بازتوانی قلبی ۱۶۵۱۸ بیمار، برای درمان با استاتین با دوز بالا ۲۳۲۴۰ بیمار، و برای درمان دوگانه ضد پلاکتی ۲۲۳۳۱ بیمار تعیین گردید. یافته‌ها نشان داد که بخش قابل توجهی از بیماران از رعایت کامل توصیه‌های پیشگیری ثانویه بازمی‌مانند؛ تنها ۱۳.۴ درصد از بیماران به برنامه‌های بازتوانی قلبی پایبند بودند که همین میزان مشارکت با کاهش معنی‌دار خطر مرگ و میر ارتباط مستقیم داشت (\lr{HR=0.73, p=0.036}). همچنین مشخص شد ۶۶.۴ درصد بیماران استاتین با دوز بالا را رها کرده‌اند که این عدم پایبندی به‌طور مستقل ریسک مرگ و میر را افزایش داد (\lr{HR=1.16, p<0.001})، و بیش از نیمی از این موارد قطع دارو در همان ماه اول رخ داده بود. پایبندی به دستورالعمل‌های درمان ضد پلاکتی نیز تنها ۲۹.۳ درصد بود، اما رعایت آن ریسک مرگ و میر را ۲۱ درصد کاهش داد (\lr{HR=0.79, p<0.001}). در بحث، نویسندگان به این نکته اشاره کردند که شکاف عمیق بین توصیه‌های بالینی و رفتار واقعی بیماران همچنان یک چالش جدی در سیستم‌های سلامت محسوب می‌شود. در نهایت، نتیجه‌گیری شد که برای بهبود پیامدهای بیماران، اتکا به سیستم‌های ارجاع داوطلبانه و برطرف کردن صرفِ موانع مالی کافی نیست، بلکه مداخلات چندبعدی، ارجاع ساختاریافته و پیگیری‌های مستمر مبتنی بر سیستم سلامت جهت افزایش پایبندی ضروری است \cite{Hammer2024Adherence}.

\item موهانداس و همکاران در یک مطالعه مورد-شاهدی در سال ۲۰۲۵ میلادی، پیش‌بینی‌کننده‌های مستقل انفارکتوس مجدد میوکارد و عدم پیروی درمانی از راهبردهای پیشگیری ثانویه را در بیماران مبتلا به بیماری عروق کرونر در منطقه روستایی آلووا در جنوب هند بررسی کردند. این مطالعه بر روی ۶۲۳ بیمار از کوهورت \lr{ENDIRA} انجام پذیرفت که شامل ۱۸۷ بیمار در گروه مورد (مبتلا به انفارکتوس مجدد) و ۴۳۶ بیمار در گروه شاهد (بدون سابقه عود) بود. پیامدهای مورد نظر در دامنه‌های مختلف اقدامات غیردارویی ارزیابی شدند. نتایج نشان داد که پیروی درمانی در دامنه فعالیت بدنی در گروه مورد ۶۷.۴ درصد و در گروه شاهد ۶۴.۹ درصد بود؛ ضمن اینکه سابقه جراحی بای‌پس عروق کرونر با کاهش ۸۰ درصدی احتمال عدم پیروی در این زمینه همراه بود (\lr{AOR=0.20, 95\% CI: 0.09-0.43}). در شاخص توده بدنی، تنها ۵۱.۹ درصد از گروه مورد و ۴۸.۹ درصد از شاهدان در محدوده طبیعی قرار داشتند و نارضایتی از خدمات بهداشتی خطر عدم پیروی از این شاخص را به شدت افزایش می‌داد (\lr{AOR=9.83, 95\% CI: 1.01-78.21}). کنترل فشار خون در ۶۶.۳ درصد از موارد و ۷۷.۳ درصد از شاهدان مطلوب بود؛ عواملی نظیر داشتن بیش از چهار عضو خانواده (\lr{AOR=2.32, 95\% CI: 1.15-4.68}) و ابتلا به بیماری بیش از ۵ سال (\lr{AOR=3.29, 95\% CI: 1.61-6.72}) با عدم پیروی مناسب در این بخش ارتباط مستقیم داشتند. همچنین، ترک کامل سیگار در ۶۰.۴ درصد از گروه مورد و ۶۱.۸ درصد از گروه شاهد گزارش شد و سلامت روان نیز در ۵۹.۴ درصد از بیماران با سابقه عود در مقایسه با ۷۲.۵ درصد از گروه شاهد در وضعیت مطلوبی قرار داشت. ارزیابی پیش‌بینی‌کننده‌های مستقل عود بیماری نشان داد که وضعیت اقتصادی-اجتماعی ضعیف (\lr{AOR=1.98, 95\% CI: 1.34-2.91})، مصرف الکل (\lr{AOR=1.65, 95\% CI: 1.04-2.61}) و فاصله بیش از ۲۰ کیلومتر تا مرکز درمانی (\lr{AOR=1.85, 95\% CI: 1.09-3.12}) از عوامل خطر کلیدی محسوب می‌شوند. محققان نتیجه گرفتند که علل انفارکتوس مجدد و عدم پایبندی به اقدامات پیشگیرانه ماهیتی چندعاملی داشته و کاهش میزان عود نیازمند تقویت دسترسی عادلانه به خدمات مراقبتی و ارتقای حمایت‌های اجتماعی در جوامع روستایی است.

\item سالومون و همکاران در مطالعه‌ای گذشته‌نگر در سال 2020 به بررسی ارتباط پایبندی تجمعی (\lr{Cumulative Adherence}) به دستورالعمل‌های پیشگیری ثانویه و میزان مرگ‌ومیر درازمدت پس از انفارکتوس حاد میوکارد پرداختند. این مطالعه کوهورت بر روی 25778 بیمار در بازه زمانی 30 روزه و 24200 بیمار در بازه 90 روزه پس از ترخیص در یک سیستم سلامت یکپارچه در کالیفرنیا انجام شد. محققان وضعیت پایبندی بیماران را به 6 تا 7 معیار اساسی شامل درمان‌های دارویی، کنترل عوامل خطر بالینی (مانند فشار خون و کلسترول) و همچنین مداخلات غیردارویی مرتبط با سبک زندگی (مانند عدم مصرف دخانیات) ارزیابی کردند. نتایج نشان داد که پایبندی همزمان به چندین دستورالعمل، به طور مستقل با کاهش چشمگیر مرگ‌ومیر همراه است؛ به گونه‌ای که در مدل ارزیابی 90 روزه، بیمارانی که به تمامی 7 دستورالعمل بهداشتی پایبند بودند، در مقایسه با گروهی که تنها 0 تا 3 معیار را رعایت می‌کردند، خطر مرگ‌ومیر بسیار کمتری داشتند (نسبت خطر \lr{[HR]} برابر با 0.57؛ فاصله اطمینان 95\%: 0.49 تا 0.66). علاوه بر این، پژوهشگران دریافتند که دستیابی به هر دستورالعمل توصیه‌شده اضافی، به تنهایی و به صورت معناداری باعث کاهش نرخ مرگ‌ومیر می‌شود (نسبت خطر 0.92؛ فاصله اطمینان 95\%: 0.90 تا 0.94). این مطالعه نتیجه‌گیری کرد که اگرچه رعایت تک‌تک توصیه‌ها فواید نسبی به همراه دارد، اما بیشترین میزان بقای بیماران تنها در سایه پایبندی همه‌جانبه به ترکیبی از دارودرمانی، کنترل عوامل خطر و تغییرات غیردارویی سبک زندگی به دست می‌آید. این یافته‌ها به روشنی خلاء تمرکز صرف بر درمان‌های دارویی را برجسته می‌سازد و نشان می‌دهد که برای ارتقای کیفیت مراقبت و کاهش مرگ‌ومیر پس از سکته قلبی، ارزیابی دقیق میزان رعایت تدابیر غیردارویی پیشگیری ثانویه توسط بیماران و شناخت عوامل بازدارنده یا تسهیل‌کننده مرتبط با آن، امری اجتناب‌ناپذیر و حیاتی است.
\end{itemize}